% !TeX root = main.tex

\section{Background}
\label{sec:Background}

QMCPACK is an open-source Quantum Monte Carlo framework with a sizeable number of users as evidenced by the number of
citations of the reference paper by Kim et al.~\cite{kim2018qmcpack}, 223 as of April 2023 according to Google Scholar.
Development
started as early as 1998 and has involved many stages of development and modernization through the years from different teams of developers~\cite{doi:10.1063/5.0004860}. QMCPACK's current development process ensures that the code performs on several heterogeneous architectures targeting a wide range of parallel computing systems. The latest efforts include its readiness on novel Arm-powered test bed clusters~\cite{10.1145/3581576.3581621}.

Table~\ref{tab:software} presents some of the characteristics of the QMCPACK software infrastructure:

\begin{table}[!ht]
\centering
\caption{Software characteristics of QMCPACK}
\begin{tabular}{| p{0.35\linewidth} | p{0.55\linewidth}|}
\hline
Software characteristic  & Value \\
\hline
Repository           & \url{https://github.com/QMCPACK/qmcpack} \\ 
Open-source License  & UIUC / NCSA \\
Languages            & C\texttt{++}17, C, Python 3 \\
Parallel models      & CUDA, HIP, OpenMP, SYCL, MPI \\
Build system         & CMake $\geq$ 3.17 \\ 
Lines of code        &  $\sim 2\times10^5$       \\
Tests suites         & unit, deterministic, stochastic, performance \\
Dependencies         & Libxml2, Boost, HDF5, BLAS, LAPACK, FFTW3, MPI \\ 
Current version      & 3.16.0 \\
\hline
\end{tabular}
\label{tab:software}
\end{table}

Compilation of the QMCPACK C\texttt{++} source code generates a binary executable per selected configuration, \texttt{qmcpack} (real number calculations) and
\texttt{qmcpack\_complex} (complex number calculations). Users can configure and build the code via CMake, but the preferred installation route is through the 
Spack~\cite{gamblin2015spack} package manager designed specifically for supercomputers and easy installation for multiple configurations on multiple operating systems. QMCPACK ships with the Nexus workflow management system~\cite{KROGEL2016154} to ease the running of research workflows in the
creation of inputs and data analysis of outputs. Nexus acts as a front end to QMCPACK, both by abstracting the interface to the application via Python calls and facilitating running QMCPACK at scale on production HPC systems.

QMCPACK's current software process can be summarized as follows: 

\begin{itemize}
    \item GitHub ``fork and pull request'' contribution workflow.
    \item Continuous integration (CI) checks on many build combinations are enforced.
    \item Code contributions must be reviewed by the core team, CI checks passed, and the final merge must be done by someone at an institution not involved with the contribution.
    \item Open source compiler and library dependencies are supported for versions released in the last two years.
    \item Security: two-factor authentication for all members and manually triggered self-hosted CI by a few admins.
    \item User documentation hosted on \texttt{readthedocs}\footnote{\url{https://qmcpack.readthedocs.io/}}.
    \item Tests that take too long to run practically in CI are run nightly on several systems and results reported on cdash\footnote{\url{https://cdash.qmcpack.org/}}.
    \item Communication is done via project Slack, GitHub issues, and regular project meetings.
\end{itemize}

QMCPACK is heavily developed and maintained by National Laboratory staff members, as indicated by the authors' affiliations. Much of the work is carried out by the core team of domain scientists in QMC methods, in collaboration with computer scientists in HPC, and those identified as RSEs at National Laboratories~\cite{10071971}. While there is no formal RSE career track, the addition of staff dedicated to software engineering tasks from Oak Ridge National Laboratory~\cite{10078175} has been welcomed in executing the tasks described in the present paper towards meeting QMCPACK's sustainability goals.
